

Supposons $q=1$. Alors $\operatorname{dim{.}supp}(N_q)=0$ ; donc $H^1(N_q(-1))=0$ et $h^0(F_q(-1)\le h^0(F(-1))\le b_0$. De plus, $G'=0$ ; donc $G\simeq G''$. Ainsi, $F_1$ est bien un $(b_0,\ldots,b_r)$-faisceau. En outre, 6.5 (ii) montre que $H^0(F_1(-b_0))=0$. Si $q\ge2$, on peut supposer $G_{q-1}$ un $((b_{q-1})^{q-1},\ldots,b_{q-1},\ldots,b_r)$-faisceau et $H^0(G_{q-1}(-b_{q-1}))=0$, en raisonnant par r\'ecurrence sur $q$. Alors 6.5 montre que $H^0(F_q(-b_{q-1}))=0$ et $h^0(F_q(-1))\le (b_{q-1})^q$ ; d'o\`u l'assertion.

\textbf{Th\'eor\`eme 6.7.} \emph{Il existe deux suites $\{A_i(X_0,\ldots,X_i;Y)\}$ et $\{A_i^{(q)}(X_0,\ldots,X_q;Y)\}$ de polyn\^omes ayant les propri\'et\'es suivantes. On suppose les conditions de 6.1 remplies et $F$ un quotient de ${O}_P(-m)^{\oplus M}$ avec $m\ge0$.}
\begin{enumerate}[label=(\roman*)]
\item \emph{Si $F$ est un $(b_0,\ldots,b_r)$-faisceau, alors pour $i=0,\ldots,r$,}
\[
   |a_i|\le A_{r-i}(b_i,\ldots,b_r;m).
\]
\item \emph{Si $\chi(F_q(n))=\sum_{i=0}^{r} a_i^{(q)}\binom{n+i}{i}$ est le polyn\^ome de Hilbert de $F_q$, alors pour $i=0,\ldots,q-1$,}
\[
   |a_i^{(q)}|\le A_{r-i}^{(r-q)}(c_{q-1},\ldots,c_r;m).
\]
\end{enumerate}
\emph{(Remarquons que $a_{q-1}^{(q)}\le a_{q-1}$ et $a_q^{(q)}$ et $a_q^{(q)}=a_q,\ldots,a_r^{(q)}=a_r$ puisque $\operatorname{dim{.}supp}(N_q)\le q-1$).}

En effet, soit une suite exacte $0\to F(-1\to F\to G\to0$ telle que $G$ soit un $(b_1,\ldots,b_r)$-faisceau ; raisonnant par r\'ecurrence sur $r$, on peut supposer d\'efinis $A_0,\ldots,A_{r-1}$, compte tenu de 1.7. Alors 6.4 (ii) met en \'evidence la d\'efinition de $A_r$, compte tenu de 1.8. Enfin (ii) r\'esulte formellement de 6.4 (iii), 6.6, et l'assertion (i) (o\`u on prend $F_q$ pour $F$).

